\section{Prior ECC2K-130 implementations}
\label{sec:history}

The 2009 attack is the baseline every number in this paper should be
read against. Bailey et al.\ fixed the iteration function, the
distinguished-point rule and the $2^{60.9}$ budget, and measured one
step on every platform they could get: Core~2 and Phenom CPUs, the
Cell processor in a PlayStation~3, the NVIDIA GTX~295, Spartan-3
FPGAs, and a synthesised ASIC~\cite{bailey2009}. Follow-up papers
then tuned the Cell~\cite{bos2010cell}, the GPU~\cite{bernstein2010gpu}
and the FPGA~\cite{fan2010fpga} clients, and Niederhagen's thesis
collects the final figures~\cite{niederhagen2012thesis}.
Table~\ref{tab:history} restates them in this paper's unit.

\begin{table}[t]
\centering
\footnotesize
\setlength{\tabcolsep}{3pt}
\caption{ECC2K-130 rates reported in 2009--2012, all for the
  Bailey et al.\ iteration function, converted to billions of
  iterations per second per chip. ``Reported'' is the source's own
  figure; the converted value is ours where the source gives cycles
  per iteration (cores $\times$ clock / cycles). GTX~295 figures are
  per card, which carries two GPU dies. The last column is the number
  of chips the source says (or its rate implies) would perform
  $2^{60.9}$ iterations in one year.}
\label{tab:history}
\begin{tabular}{@{}lllrrr@{}}
\toprule
year & platform & reported & billion it/s & chips, 1 yr & source \\
\midrule
2009 & Core 2 Q6850, 4 cores, 3.0~GHz & 534 cycles/it per core & 0.0225 & 3{,}039 & \cite{bailey2009} \\
2009 & Phenom 9550, 4 cores, 2.2~GHz & 637 cycles/it per core & 0.0138 & --- & \cite{bailey2009} \\
2009 & Cell (PS3, 6 SPEs + PPE) & 27.67 M it/s & 0.0277 & 2{,}466 & \cite{bailey2009} \\
2009 & GTX 295 (two GT200b dies) & 25.12 M it/s & 0.0251 & 2{,}716 & \cite{bailey2009} \\
2009 & Spartan-3 XC3S5000, 20 engines & 33.67 M it/s & 0.0337 & 2{,}026 & \cite{bailey2009} \\
2009 & ASIC, 90~nm (synthesis est.) & 800 M it/s & 0.800 & 85 & \cite{bailey2009} \\
2010 & Spartan-3 XC3S5000 & 111 M it/s & 0.111 & --- & \cite{fan2010fpga} \\
2010 & GTX 295, final kernel & 1{,}164 cycles/it per MP & 0.064 & 1{,}066 & \cite{bernstein2010gpu,niederhagen2012thesis} \\
2012 & Cell (PS3, 6 SPEs) & 741.79 cycles/it per SPE & 0.0259 & 2{,}636 & \cite{niederhagen2012thesis} \\
\bottomrule
\end{tabular}
\end{table}

A few details matter when comparing. The GTX~295 is two GPUs on one
card, so its $25.12$ and $64$~million iterations per second are per
card; per die they are half. The INDOCRYPT paper prints ``63 million''
for the final kernel, while its own cycle count gives $64.0$, which is
also what the thesis and the two-year device count
use~\cite{bernstein2010gpu,niederhagen2012thesis}. The 2009 FPGA figure
was partly estimated (only the multiplier of the chosen design was
implemented), the ASIC is a synthesis result rather than silicon, and
the $111$~million figure for the 2010 FPGA design is as cited by later
work, since we could not obtain the primary text. The 85-ASIC count
comes from the project's status page, which Bailey et al.'s ``less than
100 dedicated ASICs'' matches.

\begin{figure}[t]
\centering
\begin{tikzpicture}
\begin{axis}[
  xbar, xmode=log, log basis x=10, log origin=infty,
  width=0.64\textwidth, height=0.6\textwidth,
  bar width=7pt, bar shift=0pt,
  xmin=0.008, xmax=80,
  xtick={0.01,0.1,1,10}, xticklabels={0.01,0.1,1,10},
  xlabel={billion iterations per second per chip or card (log scale)},
  symbolic y coords={core2,cell09,gtx09,sthreea,cell12,gtx10,sthreeb,asic,xeon,vu47p,pro6000c,pro6000b},
  ytick={core2,cell09,gtx09,sthreea,cell12,gtx10,sthreeb,asic,xeon,vu47p,pro6000c,pro6000b},
  yticklabels={
    {2009 Core 2 Q6850 (4 cores)},
    {2009 Cell, PlayStation 3},
    {2009 GTX 295 (card, 2 dies)},
    {2009 Spartan-3 XC3S5000},
    {2012 Cell, PlayStation 3},
    {2010 GTX 295, final (card, 2 dies)},
    {2010 Spartan-3 XC3S5000},
    {2009 ASIC 90\,nm (estimate)},
    {2026 Xeon, 4 threads, AVX-512},
    {2026 VU47P FPGA (AWS F2)},
    {2026 RTX PRO 6000, campaign},
    {2026 RTX PRO 6000, best $\sigma$},
  },
  y dir=reverse,
  xmajorgrids,
  enlarge y limits=0.05,
  nodes near coords,
  nodes near coords style={font=\scriptsize,color=black,anchor=west},
  point meta=explicit symbolic,
  legend style={at={(0.5,1.02)},anchor=south},
  legend columns=3,
  legend cell align=left,
]
\addplot[fill=inkS!55,draw=none] coordinates {
  (0.0225,core2) [0.0225] (0.0277,cell09) [0.0277] (0.0251,gtx09) [0.0251]
  (0.0337,sthreea) [0.0337] (0.0259,cell12) [0.0259] (0.064,gtx10) [0.064]
  (0.111,sthreeb) [0.111]};
\addlegendentry{2009--2012, measured}
\addplot[fill=white,draw=inkS,pattern=north east lines,pattern color=inkS!70] coordinates {(0.8,asic) [0.8 est.]};
\addlegendentry{2009 synthesis estimate}
\addplot[fill=cA,draw=none] coordinates {
  (0.040,xeon) [0.040] (8.25,vu47p) [8.25] (14.64,pro6000c) [14.64] (15.88,pro6000b) [15.88]};
\addlegendentry{this paper, measured}
\end{axis}
\end{tikzpicture}
\caption{ECC2K-130 rates, 2009--2012 against today, on one log scale
  (Tables~\ref{tab:history} and~\ref{tab:history-compare}). Each bar is
  one chip except the two GTX~295 bars, which are one card carrying two
  GPU dies, as their sources report them; per die they are half. One RTX~PRO~6000
  does the work of about 229 of the 2010 cards.}
\label{fig:history}
\end{figure}


\subsection{Comparison with current hardware}
\label{sec:history-compare}

Table~\ref{tab:history-compare} sets the best 2009--2012 figure for
each platform class next to what we measured. All rows walk the same
classes of the same curve with the same $2^{60.9}$ budget; the 2009
walk and our campaign walk are the same map, and only the arithmetic
and the hardware differ.

\begin{table}[t]
\centering
\footnotesize
\caption{Best published rate per platform class against this paper's
  measurements, billions of iterations per second per chip. Ratios
  are ours. One year of $2^{60.9}$ iterations needs about $68.2$ billion
  iterations per second.}
\label{tab:history-compare}
\begin{tabular}{@{}llrlrr@{}}
\toprule
class & then & billion it/s & now & billion it/s & ratio \\
\midrule
GPU, per card & GTX 295 (2010) & 0.064 & RTX PRO 6000, campaign & 14.64 & $229\times$ \\
 & & & RTX PRO 6000, best $\sigma$ & 15.88 & $248\times$ \\
GPU, per die & GT200b (2010) & 0.032 & RTX PRO 6000 & 14.64 & $458\times$ \\
FPGA & Spartan-3 XC3S5000 (2010) & 0.111 & VU47P (AWS F2) & 8.25 & $74\times$ \\
ASIC & 90~nm estimate (2009) & 0.800 & --- & --- & --- \\
CPU, per core & Core 2, 3.0~GHz (2009) & 0.0056 & Xeon 2.1~GHz, AVX-512 & 0.0126 & $2.2\times$ \\
\bottomrule
\end{tabular}
\end{table}

The GPU row is the headline. One RTX~PRO~6000 running the campaign
build does the work of about 229 GTX~295 cards, and at that rate the
whole $2^{60.9}$ budget is about $4.7$ GPU-years instead of the 1{,}066
card-years of 2012. Of that factor, part is simply a bigger and faster
chip: 188 SMs at about $2.4$~GHz against 60 multiprocessors at
$1.24$~GHz is about $6\times$ more SM-clocks per second. The rest, roughly
$38\times$ per SM-clock, is the arithmetic: a native carry-less product
instead of a bitsliced one, three-input logic, a packed representation
with the Frobenius as a permutation network, and a batch shape tuned
against measured pipe rates. Figure~\ref{fig:gains} splits that factor, and Sections~\ref{sec:arith}
through~\ref{sec:gpu} take it apart step by step.

The FPGA gain is smaller, about $74\times$ from a 2010 Spartan-3 to a
current VU47P, and on cost per solve the F2 is level with the GPU
(Section~\ref{sec:fpga-vs-gpu}). The CPU gain per core is only about
$2\times$: the CPU client is the same bitsliced design as in 2009, and
it is the GPU that changed. The one platform class we have not
revisited is the ASIC. A 2009 synthesis estimate of $0.8$ billion
iterations per second per chip at 90~nm already needed fewer than a
hundred chips for a one-year solve; a modern node would need far
fewer, but that is a manufacturing question and outside this paper.
